% ============================================================
% LUNARA Capstone Project — Milestone 4: Development Phase (Coding)
% CST-452 Senior Capstone — Grand Canyon University
% Authors: Owen Lindsey, Andrew Mack, Carter Wright
% Compile: xelatex (requires MiKTeX with fontspec)
% ============================================================
\documentclass[12pt,letterpaper]{article}

% ── Page Geometry ──
\usepackage[top=1in, bottom=1in, left=1in, right=1in]{geometry}

% ── Typography ──
\usepackage{fontspec}
\setmainfont{Times New Roman}
\setsansfont{Arial}
\setmonofont{Consolas}[Scale=0.85]
\usepackage{setspace}
\onehalfspacing
\usepackage{parskip}
\setlength{\parindent}{0pt}

% ── Color & Graphics ──
\usepackage[dvipsnames,svgnames,x11names]{xcolor}
\definecolor{GCUNavy}{HTML}{2D3748}
\definecolor{LunaraBlue}{HTML}{4A90D9}
\definecolor{CodeBG}{HTML}{F7F8FA}
\definecolor{TableHeader}{HTML}{2D3748}
\definecolor{TableHeaderText}{HTML}{FFFFFF}
\definecolor{TableRowAlt}{HTML}{F0F4F8}
\usepackage{graphicx}
\usepackage{float}
\usepackage{tikz}

% ── Tables ──
\usepackage{array}
\usepackage{booktabs}
\usepackage{tabularx}
\usepackage{longtable}
\usepackage{multirow}
\usepackage{colortbl}
\renewcommand{\arraystretch}{1.3}

% ── Longtable page-break settings ──
\setlength{\LTpre}{0.8em}
\setlength{\LTpost}{0.8em}

% ── Code Listings ──
\usepackage{listings}
\lstdefinelanguage{TypeScript}{
  keywords={import, from, export, default, const, let, var, function, return, if, else, for, of, in, new, try, catch, async, await, class, extends, implements, interface, type, enum, as, typeof, instanceof, void, null, undefined, true, false, this, super, switch, case, break, continue, throw, while, do, with, yield, static, public, private, protected, readonly, abstract, get, set},
  keywordstyle=\color{blue}\bfseries,
  ndkeywords={string, number, boolean, any, never, unknown, Promise, Array, Record, Date, Document, Schema, Express, Request, Response, NextFunction},
  ndkeywordstyle=\color{teal}\bfseries,
  sensitive=true,
  comment=[l]{//},
  morecomment=[s]{/*}{*/},
  commentstyle=\color{gray}\itshape,
  stringstyle=\color{orange!80!black},
  string=[b]',
  morestring=[b]",
  morestring=[b]`,
}
\lstset{
  language=TypeScript,
  basicstyle=\ttfamily\scriptsize,
  backgroundcolor=\color{CodeBG},
  frame=single,
  framerule=0.5pt,
  rulecolor=\color{gray!40},
  breaklines=true,
  breakatwhitespace=false,
  showstringspaces=false,
  tabsize=2,
  numberstyle=\tiny\color{gray},
  numbers=left,
  numbersep=8pt,
  xleftmargin=1.5em,
  framexleftmargin=1.5em,
  aboveskip=1em,
  belowskip=1em,
  captionpos=b,
  escapeinside={(*@}{@*)},
  literate={=>}{{\textcolor{blue}{=>}}}{2}
           {===}{{\textcolor{blue}{===}}}{3}
           {!==}{{\textcolor{blue}{!==}}}{3}
           {??}{{\textcolor{blue}{??}}}{2},
}

% ── Hyperlinks ──
\usepackage[colorlinks=true, linkcolor=GCUNavy, urlcolor=LunaraBlue, citecolor=GCUNavy, bookmarks=true, bookmarksnumbered=true]{hyperref}
\urlstyle{same}

% ── Headers & Footers ──
\usepackage{fancyhdr}
\setlength{\headheight}{14pt}
\addtolength{\topmargin}{-2pt}
\pagestyle{fancy}
\fancyhf{}
\fancyhead[L]{\small\textcolor{GCUNavy}{\textsc{LUNARA Capstone Project}}}
\fancyhead[R]{\small\textcolor{GCUNavy}{\textsc{Milestone 4 --- Coding}}}
\fancyfoot[L]{\small\textcolor{gray}{Grand Canyon University}}
\fancyfoot[C]{\small\textcolor{gray}{CST-452}}
\fancyfoot[R]{\small\textcolor{gray}{\thepage}}
\renewcommand{\headrulewidth}{0.4pt}
\renewcommand{\footrulewidth}{0.2pt}

% ── Section Formatting ──
\usepackage{titlesec}
\titleformat{\section}{\Large\bfseries\color{GCUNavy}}{\thesection}{1em}{}[\vspace{0.3em}{\color{LunaraBlue!40}\rule{\textwidth}{0.75pt}}\vspace{0.2em}]
\titleformat{\subsection}{\large\bfseries\color{GCUNavy}}{\thesubsection}{1em}{}
\titleformat{\subsubsection}{\normalsize\bfseries\color{GCUNavy}}{\thesubsubsection}{1em}{}

% ── Math ──
\usepackage{amsmath}

% ── Captions ──
\usepackage{caption}
\captionsetup{font=small, labelfont={bf,color=GCUNavy}, textfont={color=black}, skip=8pt}

% ── Enumerate / Itemize ──
\usepackage{enumitem}

% ── Page-break helpers ──
\usepackage{needspace}

% ── PDF Metadata ──
\hypersetup{
  pdftitle={LUNARA Postpartum Support Platform --- Milestone 4: Development Phase (Coding)},
  pdfauthor={Owen Lindsey, Andrew Mack, Carter Wright},
  pdfsubject={CST-452 Senior Capstone --- Grand Canyon University},
  pdfkeywords={capstone, postpartum, LUNARA, full-stack, React, Node.js, MongoDB},
}

% ── Custom Commands ──
\newcommand{\screenshotfig}[3]{%
  \begin{figure}[H]
    \centering
    \includegraphics[width=0.85\textwidth,height=0.75\textheight,keepaspectratio]{#1}
    \caption{#2}
    \label{fig:#3}
  \end{figure}
}

% ============================================================
\begin{document}

% ── TITLE PAGE ──
\thispagestyle{empty}
\begin{tikzpicture}[remember picture, overlay]
  % Navy background
  \fill[GCUNavy] (current page.south west) rectangle (current page.north east);
  % Subtle accent line below title block
  \draw[LunaraBlue!60, line width=2pt] ([xshift=1.5in, yshift=-14.5cm]current page.north west) -- ([xshift=-1.5in, yshift=-14.5cm]current page.north east);
\end{tikzpicture}

\begin{center}
\vspace*{2.5cm}

{\color{white}\fontsize{32}{38}\selectfont\textbf{LUNARA}}

\vspace{0.4cm}

{\color{LunaraBlue!70}\fontsize{14}{18}\selectfont\textsc{Postpartum Support Platform}}

\vspace{1.2cm}

{\color{white}\Large Milestone 4: Development Phase --- Coding}

\vspace{0.3cm}

{\color{white!80}\large Capstone Project Paper Packet}

\vspace{3cm}

{\color{white}\large\textbf{Owen Lindsey}}

\vspace{0.15cm}

{\color{white!80}\normalsize Development Lead}

\vspace{0.8cm}

{\color{white}\large\textbf{Andrew Mack}}

\vspace{0.15cm}

{\color{white!80}\normalsize Backend Developer}

\vspace{0.8cm}

{\color{white}\large\textbf{Carter Wright}}

\vspace{0.15cm}

{\color{white!80}\normalsize Frontend Developer}

\vspace{2cm}

{\color{white!90}\normalsize Grand Canyon University}

{\color{white!90}\normalsize College of Science, Engineering and Technology}

{\color{white!90}\normalsize CST-452: Senior Capstone Project II}

{\color{white!90}\normalsize Professor Amr Elchouemi}

\vspace{0.5cm}

{\color{LunaraBlue}\normalsize March 10, 2026}

\vspace{0.6cm}

{\color{LunaraBlue!80}\normalsize \url{https://www.lunaracare.org/}}

\vspace{0.3cm}

{\color{LunaraBlue!80}\normalsize Video Demo: \url{https://www.youtube.com/watch?v=EY6ncwOMFUQ}}

\end{center}

\newpage

% ── TABLE OF CONTENTS ──
\thispagestyle{fancy}
\tableofcontents
\newpage

% ============================================================
% SECTION 1: CODING OBJECTIVE
% ============================================================
\section{Coding Objective}

The main objective of the development phase is to build the project code and create documentation based on the requirements and analysis decisions. This phase transforms the design documents into an executable program. System entities such as objects, data tables, classes, and components are developed, integrated, and tested to create a working application. The coding objective for Milestone~4 was to transform the approved requirements from the CST-451 Project Requirements document (FR1--FR15, NFR1--NFR10) and Architecture Plan into integrated, executable software modules while preserving scope discipline, quality, and maintainability. This milestone represented the most intensive period of feature construction in the LUNARA project, producing the working codebase that would later be formally validated in Milestone~5 and packaged for release in Milestone~6.

\section{Development Deliverables}

This document provides the complete submission of the Implementation Plan with detailed task breakdown, the Functional Requirements Mapping (Traceability Matrix), Module Test Cases with testing framework context, the Source Code Listing with architectural documentation, and the Application Demonstration plan. The code repository is hosted on GitHub at \url{https://github.com/omniV1/lunaraCare}.

\newpage

% ============================================================
% SECTION 2: IMPLEMENTATION PLAN
% ============================================================
\section{Implementation Plan}

\subsection{Methodology and Execution Framework}

\textbf{Formal Methodology Adopted:} Agile Scrum with iterative coding loops\\
\textbf{Planning Tools:} GitHub Project Boards, sprint planning documents in \texttt{Docs/Planning/SPRINT\_PLAN.md}\\
\textbf{Code Review Process:} GitHub Pull Request workflow with peer review\\
\textbf{Testing Integration:} Continuous integration with automated test execution

Each iteration cycle followed a consistent pattern: define a scoped story objective, implement the backend route/service/model path, build the frontend consumption layer and UI behavior, integrate and validate in the running environment, and refine based on testability and usability constraints.

\subsection{Implementation Plan Matrix --- CST-452 Development Iterations}

{\small
\begin{longtable}{|>{\raggedright\arraybackslash}p{2.8cm}|>{\raggedright\arraybackslash}p{6.5cm}|c|c|c|}
\caption{Implementation Plan Matrix --- CST-452 Development Iterations}\label{tab:impl-plan}\\
\hline
\rowcolor{TableHeader}\textcolor{TableHeaderText}{\textbf{Use Case / User Story}} & \textcolor{TableHeaderText}{\textbf{Detailed Development Tasks}} & \textcolor{TableHeaderText}{\textbf{Est.}} & \textcolor{TableHeaderText}{\textbf{Act.}} & \textcolor{TableHeaderText}{\textbf{\%}} \\
\hline
\endfirsthead
\multicolumn{5}{l}{\small\itshape\color{gray} Table~\ref{tab:impl-plan} continued from previous page}\\[0.5em]
\hline
\rowcolor{TableHeader}\textcolor{TableHeaderText}{\textbf{Use Case / User Story}} & \textcolor{TableHeaderText}{\textbf{Detailed Development Tasks}} & \textcolor{TableHeaderText}{\textbf{Est.}} & \textcolor{TableHeaderText}{\textbf{Act.}} & \textcolor{TableHeaderText}{\textbf{\%}} \\
\hline
\endhead
\hline
\multicolumn{5}{r}{\small\itshape\color{gray} Continued on next page}\\
\endfoot
\hline
\endlastfoot

\textbf{FR1: Public Website} & Landing page with hero, offerings accordion, inquiry form; contact POST endpoint & 12 & 12 & 100\% \\
\hline
\rowcolor{TableRowAlt}
\textbf{FR2: Authentication \& Registration} & JWT + Passport.js local/JWT strategies; register, login, email verify, password reset, token refresh, MFA (TOTP) & 24 & 26 & 100\% \\
\hline
\textbf{FR3: User Dashboards} & Provider dashboard (overview, clients, schedule, blog, profile, createProvider tabs); Client dashboard (overview, messages, appointments, profile tabs) with embedded check-ins, documents, resources, intake, and care plan views & 32 & 34 & 100\% \\
\hline
\rowcolor{TableRowAlt}
\textbf{FR4: Dynamic Intake Forms} & Five-step wizard (personal, birth, feeding, support, health); Zod validation per step; auto-save; section-level PATCH & 16 & 18 & 100\% \\
\hline
\textbf{FR5: Real-time Messaging} & Socket.IO server with JWT auth; user rooms; conversation rooms; send/receive/deliver events; message rate limiter; MongoDB persistence; MessageCenter and ClientMessageProvider components & 20 & 22 & 100\% \\
\hline
\rowcolor{TableRowAlt}
\textbf{FR6: Appointment Scheduling} & CRUD endpoints; upcoming/request/confirm/cancel flows; appointment notification service; reminder service; provider calendar and client scheduling modal & 16 & 16 & 100\% \\
\hline
\textbf{FR7: Resource Library} & Resource CRUD with versioning; category management; client-side browsable library with detail modal; resource interactions (favorites, stats) & 12 & 12 & 100\% \\
\hline
\rowcolor{TableRowAlt}
\textbf{FR8: Daily Check-ins} & Check-in model (mood 1--10, 10 symptoms, notes, share toggle); trend service; alert service; provider review panel; client check-in form with history and trend visualization & 16 & 18 & 100\% \\
\hline
\textbf{FR9: Provider Client Management} & Client listing with assign/unassign; client creation with auto-generated password; provider profile editing; reports and analytics & 12 & 14 & 100\% \\
\hline
\rowcolor{TableRowAlt}
\textbf{FR10: Care Plan System} & Care plan templates with sections and milestones; plan creation and CRUD; milestone status tracking; auto-progress calculation; CarePlanManager component & 16 & 16 & 100\% \\
\hline
\textbf{FR11: Blog Platform} & Rich text editor; draft/publish workflow; blog versioning; public blog list and detail pages; blog management interface & 10 & 10 & 100\% \\
\hline
\rowcolor{TableRowAlt}
\textbf{Documents \& Files} & Document CRUD with type and privacy levels; status workflow (draft\textrightarrow submitted\textrightarrow reviewed\textrightarrow completed); GridFS file storage; bulk upload; provider document review with feedback & 14 & 16 & 100\% \\
\hline
\textbf{Push Notifications} & Web Push with VAPID keys; subscription management; test notification endpoint & 6 & 6 & 100\% \\
\hline
\rowcolor{TableRowAlt}
\textbf{Security Middleware} & Helmet, CORS, rate limiting, input validation (express-validator), compression, morgan logging & 8 & 8 & 100\% \\
\hline
\textbf{Administrative} & Provider account creation; admin route module & 4 & 4 & 100\% \\
\hline
\rowcolor{TableRowAlt}
\textbf{FR12: Digital Journaling} & --- & --- & --- & Deferred \\
\hline
\textbf{FR13: Daily Insights} & --- & --- & --- & Deferred \\
\hline
\rowcolor{TableRowAlt}
\textbf{FR14: Sleep \& Feeding Trackers} & --- & --- & --- & Deferred \\
\hline
\textbf{FR15: AI Note Summarization} & --- & --- & --- & Deferred \\
\hline
\end{longtable}
}

\textbf{Percent of user stories complete for this iteration:} 100\% of in-scope stories (15/15)\\
\textbf{Percent of user stories complete for entire project:} 73\% (11 of 15 total FRs; 4 deferred to post-launch)

\begin{align}
\text{Iteration Completion Rate} &= \frac{\text{Completed Stories}}{\text{Planned Stories}} \times 100\% = \frac{15}{15} \times 100\% = 100\% \\[1em]
\text{Project Completion Rate} &= \frac{11}{15} \times 100\% \approx 73\%
\end{align}

\newpage

% ============================================================
% SECTION 3: TRACEABILITY MATRIX
% ============================================================
\section{Functional Requirements Mapping (Traceability Matrix)}

For each requirement $R_i$, at least one implementation artifact and one validation artifact are identified:
\[
R_i \Rightarrow \{M_i, T_i\}
\]
where $M_i$ represents the module(s) implementing requirement $R_i$ and $T_i$ represents the test case(s) validating requirement $R_i$.

{\scriptsize
\begin{longtable}{|>{\raggedright\arraybackslash}p{2.2cm}|>{\raggedright\arraybackslash}p{2.5cm}|>{\raggedright\arraybackslash}p{5.5cm}|>{\raggedright\arraybackslash}p{2.5cm}|}
\caption{Functional Requirements Traceability Matrix}\label{tab:trace-matrix}\\
\hline
\rowcolor{TableHeader}\textcolor{TableHeaderText}{\textbf{Functional Requirement}} & \textcolor{TableHeaderText}{\textbf{Architecture Section}} & \textcolor{TableHeaderText}{\textbf{Code Module(s)}} & \textcolor{TableHeaderText}{\textbf{Test Case(s)}} \\
\hline
\endfirsthead
\multicolumn{4}{l}{\small\itshape\color{gray} Table~\ref{tab:trace-matrix} continued from previous page}\\[0.5em]
\hline
\rowcolor{TableHeader}\textcolor{TableHeaderText}{\textbf{Functional Requirement}} & \textcolor{TableHeaderText}{\textbf{Architecture Section}} & \textcolor{TableHeaderText}{\textbf{Code Module(s)}} & \textcolor{TableHeaderText}{\textbf{Test Case(s)}} \\
\hline
\endhead
\hline
\multicolumn{4}{r}{\small\itshape\color{gray} Continued on next page}\\
\endfoot
\hline
\endlastfoot

FR1.1: Responsive landing page & Frontend Architecture & \texttt{Lunara/src/pages/LandingPage.tsx}, \texttt{Lunara/src/components/layout/MainLayout.tsx} & TC-001 \\
\hline
\rowcolor{TableRowAlt}
FR1.3: Contact/inquiry form & Form Handling & \texttt{Lunara/src/pages/LandingPage.tsx}, \texttt{backend/src/routes/public.ts} & TC-002 \\
\hline
FR2.1: Email/password registration & Authentication Service & \texttt{backend/src/routes/auth.ts}, \texttt{backend/src/config/passport.ts} & TC-003 \\
\hline
\rowcolor{TableRowAlt}
FR2.2: JWT token management & Security Architecture & \texttt{backend/src/utils/tokenUtils.ts}, \texttt{backend/src/config/passport.ts} & TC-004, TC-005 \\
\hline
FR2.3: MFA (TOTP) & Authentication & \texttt{backend/src/routes/mfa.ts} & TC-006 \\
\hline
\rowcolor{TableRowAlt}
FR3.1: Provider dashboard & Dashboard Architecture & \texttt{Lunara/src/pages/ProviderDashboard.tsx}, \texttt{Lunara/src/components/provider/*} & TC-007 \\
\hline
FR3.2: Client dashboard & Dashboard Architecture & \texttt{Lunara/src/pages/ClientDashboard.tsx}, \texttt{Lunara/src/components/client/*} & TC-008 \\
\hline
\rowcolor{TableRowAlt}
FR4.1: Dynamic intake & Intake Form & \texttt{Lunara/src/components/intake/ClientIntakeWizard.tsx}, \texttt{backend/src/routes/intake.ts} & TC-009 \\
\hline
FR5.1: Real-time messaging & Real-time Communication & \texttt{backend/src/server.ts} (Socket.IO), \texttt{Lunara/src/components/MessageCenter.tsx} & TC-010 \\
\hline
\rowcolor{TableRowAlt}
FR6.1: Appointment scheduling & Scheduling Architecture & \texttt{backend/src/routes/appointments.ts}, \texttt{Lunara/src/components/provider/ProviderCalendar.tsx} & TC-011 \\
\hline
FR7.1: Resource library & Content Management & \texttt{backend/src/routes/resources.ts}, \texttt{Lunara/src/components/resource/ResourceLibrary.tsx} & TC-012 \\
\hline
\rowcolor{TableRowAlt}
FR8.1: Daily check-ins & Wellness Tracking & \texttt{backend/src/routes/checkins.ts}, \texttt{backend/src/services/checkinTrendService.ts} & TC-013, TC-014 \\
\hline
FR9.1: Client management & Provider Dashboard & \texttt{backend/src/routes/client.ts}, \texttt{backend/src/routes/providers.ts} & TC-015 \\
\hline
\rowcolor{TableRowAlt}
FR10.1: Care plan templates & Care Plan Architecture & \texttt{backend/src/routes/carePlans.ts}, \texttt{backend/src/models/CarePlan.ts} & TC-016 \\
\hline
FR11.1: Blog publishing & Content Management & \texttt{backend/src/routes/blog.ts}, \texttt{Lunara/src/components/blog/BlogEditor.tsx} & TC-017 \\
\hline
\rowcolor{TableRowAlt}
NFR1.1: Password hashing & Security Architecture & \texttt{backend/src/models/User.ts} (bcrypt 12 rounds, pre-save hook) & TC-018 \\
\hline
NFR1.2: Input validation & Validation Architecture & \texttt{backend/src/routes/*.ts} (express-validator), Zod schemas & TC-019 \\
\hline
\rowcolor{TableRowAlt}
NFR4.1: WebSocket management & Real-time Architecture & \texttt{backend/src/services/socketConnectionManager.ts} & TC-020 \\
\hline
\end{longtable}
}

\newpage

% ============================================================
% SECTION 4: SOURCE CODE LISTING
% ============================================================
\section{Source Code Listing \& Architectural Implementation}

\subsection{System Architecture Overview}

The LUNARA platform implements a layered full-stack architecture with clear separation of concerns. The backend follows an Express.js \textbf{route \textrightarrow{} middleware \textrightarrow{} service \textrightarrow{} model} pattern, while the frontend follows a React \textbf{page \textrightarrow{} component \textrightarrow{} context \textrightarrow{} service} pattern. Both layers are written entirely in TypeScript for end-to-end type safety.

\subsection{Backend Architecture Implementation (\texttt{/backend})}

\subsubsection{Core Application Layer --- \texttt{src/server.ts}}

The Express application entry point configures the middleware stack, database connection, route registration, Socket.IO real-time messaging, and Swagger API documentation. The security-first middleware pipeline applies Helmet HTTP header hardening (with CSP intentionally disabled for JSON-only API endpoints), CORS policy enforcement with a configurable allowed-origins list, compression, and rate limiting differentiated between development (1,000 requests per 15~minutes) and production (100 requests per 15~minutes).

\begin{lstlisting}[caption={Server initialization and security middleware (\texttt{server.ts})}]
import express from 'express';
import mongoose from 'mongoose';
import cors from 'cors';
import helmet from 'helmet';
import compression from 'compression';
import rateLimit from 'express-rate-limit';
import http from 'node:http';
import { Server as SocketIOServer } from 'socket.io';
import passport from 'passport';
import swaggerJsdoc from 'swagger-jsdoc';
import swaggerUi from 'swagger-ui-express';

const app = express();
const server = http.createServer(app);
const io = new SocketIOServer(server, {
  cors: {
    origin: process.env.FRONTEND_URL ?? 'http://localhost:3000',
    methods: ['GET', 'POST'],
    credentials: true,
  },
});

// Security Middleware
app.use(helmet({ contentSecurityPolicy: false }));
app.use(compression());
app.use(cors({ origin: allowedOriginsSet, credentials: true }));

const limiter = rateLimit({
  windowMs: 15 * 60 * 1000,
  max: isProduction ? 100 : 1000,
  message: 'Too many requests from this IP, please try again later.',
});
\end{lstlisting}

Twenty route modules are registered under the \texttt{/api} prefix:

\begin{lstlisting}[caption={API route registration (\texttt{server.ts})}]
app.use('/api/auth', authRouter);
app.use('/api/auth/mfa', mfaRouter);
app.use('/api/appointments', appointmentsRouter);
app.use('/api/messages', messagesRouter);
app.use('/api/resources', resourcesRouter);
app.use('/api/categories', categoriesRouter);
app.use('/api/checkins', checkinsRouter);
app.use('/api/intake', intakeRouter);
app.use('/api/care-plans', carePlansRouter);
app.use('/api/admin', adminRouter);
app.use('/api/blog', blogRouter);
app.use('/api/users', usersRouter);
app.use('/api/providers', providersRouter);
app.use('/api/client', clientRouter);
app.use('/api/public', publicRouter);
app.use('/api/documents', documentsRouter);
app.use('/api/recommendations', recommendationsRouter);
app.use('/api/files', filesRouter);
app.use('/api/push', pushNotificationsRouter);
app.use('/api/interactions', interactionsRouter);
\end{lstlisting}

\newpage

\subsubsection{Authentication \& Authorization Layer --- \texttt{src/config/passport.ts}}

Authentication uses Passport.js with JWT, Local, and Google OAuth strategies. The JWT strategy extracts the bearer token from the Authorization header, verifies it against the JWT secret, and retrieves the corresponding user from MongoDB:

\begin{lstlisting}[caption={Passport.js authentication strategies (\texttt{passport.ts})}]
import passport from 'passport';
import { Strategy as JwtStrategy, ExtractJwt } from 'passport-jwt';
import { Strategy as LocalStrategy } from 'passport-local';
import bcrypt from 'bcryptjs';
import User from '../models/User';

// JWT Strategy
passport.use(
  new JwtStrategy(
    {
      jwtFromRequest: ExtractJwt.fromAuthHeaderAsBearerToken(),
      secretOrKey: process.env.JWT_SECRET ?? 'your-secret-key',
    },
    async (jwt_payload, done) => {
      try {
        const user = await User.findById(jwt_payload.id);
        if (user) return done(null, user);
        return done(null, false);
      } catch (error) {
        return done(error, false);
      }
    }
  )
);

// Local Strategy (Email/Password) with bcrypt comparison
passport.use(
  new LocalStrategy(
    { usernameField: 'email', passwordField: 'password' },
    async (email, password, done) => {
      const user = await User.findOne({
        email: email.toLowerCase()
      }).select('+password');
      if (!user || !user.password)
        return done(null, false, { message: 'Invalid credentials' });
      const isMatch = await bcrypt.compare(password, user.password);
      if (!isMatch)
        return done(null, false, { message: 'Invalid credentials' });
      return done(null, user);
    }
  )
);
\end{lstlisting}

Protected routes use \texttt{passport.authenticate('jwt', \{ session: false \})} as middleware, ensuring every request is verified before reaching the handler.

\newpage

\subsubsection{Data Access Layer --- Database Schema Design}

Nineteen Mongoose models define the domain. Representative schemas are shown below.

\paragraph{Users Collection --- \texttt{src/models/User.ts}}
\mbox{}\\
The User model represents all users in the system (clients, providers, admins) and stores authentication data, MFA configuration, and refresh tokens:

\begin{lstlisting}[caption={User model interface (\texttt{User.ts})}]
export interface IUser extends Document {
  firstName: string;
  lastName: string;
  email: string;
  password?: string;
  role: 'client' | 'provider' | 'admin';
  isEmailVerified: boolean;
  oauthProviders: IOAuthProvider[];
  profile: IProfile;
  mfaEnabled: boolean;
  mfaSecret?: string;
  mfaBackupCodes: string[];
  refreshTokens: IRefreshToken[];
  intakeCompleted?: boolean;
  intakeCompletedAt?: Date;
  comparePassword(candidatePassword: string): Promise<boolean>;
  canLogin(): boolean;
  getPermissions(): string[];
}
\end{lstlisting}

\paragraph{CheckIn Collection --- \texttt{src/models/CheckIn.ts}}
\mbox{}\\
The CheckIn model captures daily wellness self-assessments with mood rating, physical symptoms, and provider-sharing controls:

\begin{lstlisting}[caption={CheckIn model with symptom tracking (\texttt{CheckIn.ts})}]
export const PHYSICAL_SYMPTOMS = [
  'fatigue', 'sleep_issues', 'appetite_changes', 'anxiety', 'pain',
  'headache', 'nausea', 'dizziness', 'breast_soreness', 'bleeding',
] as const;

const checkInSchema = new Schema<ICheckInDocument>({
  userId: { type: Schema.Types.ObjectId, ref: 'User', required: true },
  date: { type: Date, required: true },
  moodScore: { type: Number, required: true, min: 1, max: 10 },
  physicalSymptoms: [{ type: String, enum: PHYSICAL_SYMPTOMS }],
  notes: { type: String, maxlength: 2000 },
  sharedWithProvider: { type: Boolean, default: false },
  providerReviewed: { type: Boolean, default: false },
}, { timestamps: true });

checkInSchema.index({ userId: 1, date: -1 }, { unique: true });
\end{lstlisting}

\paragraph{Appointments Collection --- \texttt{src/models/Appointment.ts}}

\begin{lstlisting}[caption={Appointment model (\texttt{Appointment.ts})}]
export interface IAppointmentDocument extends Document {
  clientId: mongoose.Types.ObjectId;
  providerId: mongoose.Types.ObjectId;
  startTime: Date;
  endTime: Date;
  status: 'requested' | 'confirmed' | 'scheduled'
        | 'completed' | 'cancelled';
  type?: 'virtual' | 'in_person';
  notes?: string;
  requestedBy: mongoose.Types.ObjectId;
  reminderSentAt?: Date;
}

appointmentSchema.index({ clientId: 1, startTime: -1 });
appointmentSchema.index({ providerId: 1, startTime: -1 });
appointmentSchema.index({ status: 1 });
\end{lstlisting}

\paragraph{Messages Collection --- \texttt{src/models/Message.ts}}

\begin{lstlisting}[caption={Message model (\texttt{Message.ts})}]
const messageSchema = new Schema<IMessageDocument>({
  conversationId: {
    type: Schema.Types.ObjectId, required: true, index: true
  },
  sender: { type: Schema.Types.ObjectId, ref: 'User', required: true },
  receiver: { type: Schema.Types.ObjectId, ref: 'User', required: true },
  content: { type: String, required: true },
  read: { type: Boolean, default: false },
  type: {
    type: String,
    enum: ['text', 'image', 'file', 'system'],
    default: 'text'
  },
}, { timestamps: { createdAt: true, updatedAt: false } });
\end{lstlisting}

\paragraph{CarePlan Collection --- \texttt{src/models/CarePlan.ts}}
\mbox{}\\
The CarePlan model includes a pre-save hook that automatically computes progress as the percentage of completed milestones:

\begin{lstlisting}[caption={CarePlan model with auto-progress hook (\texttt{CarePlan.ts})}]
export interface IMilestone {
  title: string;
  description?: string;
  weekOffset: number;
  category: 'physical' | 'emotional' | 'feeding'
           | 'self_care' | 'general';
  status: 'pending' | 'in_progress' | 'completed' | 'skipped';
}

const carePlanSchema = new Schema<ICarePlanDocument>({
  clientId: {
    type: Schema.Types.ObjectId, ref: 'User', required: true
  },
  providerId: {
    type: Schema.Types.ObjectId, ref: 'User', required: true
  },
  title: { type: String, required: true, maxlength: 200 },
  sections: [sectionSchema],
  status: {
    type: String,
    enum: ['active', 'completed', 'paused', 'archived'],
    default: 'active'
  },
  progress: { type: Number, default: 0, min: 0, max: 100 },
}, { timestamps: true });

carePlanSchema.pre('save', function (next) {
  const allMilestones = this.sections.flatMap(s => s.milestones);
  if (allMilestones.length === 0) this.progress = 0;
  else {
    const done = allMilestones.filter(
      m => m.status === 'completed' || m.status === 'skipped'
    ).length;
    this.progress = Math.round(
      (done / allMilestones.length) * 100
    );
  }
  next();
});
\end{lstlisting}

\paragraph{ClientDocument Collection --- \texttt{src/models/ClientDocument.ts}}
\mbox{}\\
Documents support type classification, privacy levels, and a status workflow:

\begin{lstlisting}[caption={ClientDocument model (\texttt{ClientDocument.ts})}]
export interface IClientDocument extends Document {
  title: string;
  documentType: 'emotional-survey' | 'health-assessment'
    | 'personal-assessment' | 'progress-photo' | 'other';
  uploadedBy: mongoose.Types.ObjectId;
  assignedProvider?: mongoose.Types.ObjectId;
  files: IFileAttachment[];
  submissionStatus: 'draft' | 'submitted-to-provider'
    | 'reviewed-by-provider' | 'completed';
  privacyLevel: 'client-only' | 'client-and-provider'
    | 'care-team';
}
\end{lstlisting}

\newpage

\subsubsection{API Layer \& Business Services}

\paragraph{Authentication Endpoints --- \texttt{src/routes/auth.ts}}
\mbox{}\\
The auth router exposes registration with express-validator input validation, login via Passport Local strategy, token refresh, email verification, password reset, and logout with token revocation. A rate limiter restricts sensitive endpoints to 5 requests per 15~minutes:

\begin{lstlisting}[caption={Auth route registration endpoint (\texttt{auth.ts})}]
const sensitiveEndpointLimiter = rateLimit({
  windowMs: 15 * 60 * 1000,
  max: 5,
  message: 'Too many attempts, please try again later.',
});

router.post('/register', [
  body('firstName').notEmpty(),
  body('lastName').notEmpty(),
  body('email').isEmail().normalizeEmail(),
  body('password').isLength({ min: 8 }),
  body('role').isIn(['client', 'provider']),
], async (req, res) => { /* registration logic */ });
\end{lstlisting}

\paragraph{Check-in Endpoints --- \texttt{src/routes/checkins.ts}}
\mbox{}\\
The check-in router validates mood scores (1--10), symptom values against the predefined enum, and note length, then saves and triggers alert analysis:

\begin{lstlisting}[caption={Check-in creation endpoint (\texttt{checkins.ts})}]
router.post('/', authenticate, [
  body('date').isISO8601(),
  body('moodScore').isInt({ min: 1, max: 10 }),
  body('physicalSymptoms').optional().isArray()
    .custom((arr) => arr.every(
      s => PHYSICAL_SYMPTOMS.includes(s)
    )),
  body('notes').optional().isString().isLength({ max: 2000 }),
  body('sharedWithProvider').optional().isBoolean(),
], async (req, res) => {
  const checkIn = new CheckIn({
    userId: user._id, ...req.body
  });
  await checkIn.save();
  const alerts = await getCheckInAlerts(
    user._id.toString()
  );
  return res.status(201).json({
    message: 'Check-in recorded', checkIn, alerts
  });
});
\end{lstlisting}

\paragraph{Check-in Trend Service --- \texttt{src/services/checkinTrendService.ts}}
\mbox{}\\
The trend service computes mood averages, daily mood mappings, symptom frequency distributions, mood direction (improving/stable/declining), and clinical alerts for low mood streaks, declining trends, and persistent symptoms:

\begin{lstlisting}[caption={Clinical alert thresholds and detection (\texttt{checkinTrendService.ts})}]
const ALERT_THRESHOLDS = {
  LOW_MOOD_SCORE: 3,
  LOW_MOOD_STREAK: 3,
  DECLINING_WINDOW: 7,
  PERSISTENT_SYMPTOM_DAYS: 5,
};

export async function getCheckInAlerts(
  userId: string
): Promise<AlertInfo[]> {
  const alerts: AlertInfo[] = [];
  const recent = await CheckIn.find({ userId })
    .sort({ date: -1 })
    .limit(ALERT_THRESHOLDS.DECLINING_WINDOW);

  // Low mood streak detection
  let lowStreak = 0;
  for (const c of recent) {
    if (c.moodScore <= ALERT_THRESHOLDS.LOW_MOOD_SCORE)
      lowStreak++;
    else break;
  }
  if (lowStreak >= ALERT_THRESHOLDS.LOW_MOOD_STREAK) {
    alerts.push({
      type: 'low_mood',
      message: `Mood has been ${
        ALERT_THRESHOLDS.LOW_MOOD_SCORE
      } or below for ${lowStreak} consecutive check-ins`,
      severity: lowStreak >= 5 ? 'critical' : 'warning',
    });
  }
  // ... declining trend and persistent symptom detection
}
\end{lstlisting}

\paragraph{Real-time Messaging --- Socket.IO in \texttt{src/server.ts}}
\mbox{}\\
Socket.IO connections are authenticated with the same JWT strategy. After verification, users join private rooms keyed to their user ID. Message delivery persists to MongoDB, broadcasts to the conversation room, sends a notification to the receiver's user room, and confirms delivery to the sender:

\begin{lstlisting}[caption={Socket.IO authentication and message delivery (\texttt{server.ts})}]
io.use((socket, next) => {
  const token = socket.handshake.auth?.token
    || socket.handshake.query?.token;
  if (!token)
    return next(
      new Error('Authentication error: missing token')
    );
  const payload = verifyAccessToken(token);
  socket.data.userId = payload.id;
  socket.data.role = payload.role;
  next();
});

io.on('connection', socket => {
  socket.on('send_message', async (messageData) => {
    if (isRateLimited(socket.data.userId)) {
      socket.emit('rate_limit', {
        error: 'Too many messages'
      });
      return;
    }
    const doc = new Message({
      ...messageData, read: false
    });
    await doc.save();
    io.to(doc.conversationId).emit(
      'new_message', payload
    );
    io.to(doc.receiver).emit(
      'new_message_notification', payload
    );
    socket.emit('message_delivered', {
      messageId: doc._id
    });
  });
});
\end{lstlisting}

\newpage

\subsection{Frontend Architecture Implementation (\texttt{/Lunara})}

\subsubsection{Application Root --- \texttt{src/App.tsx}}

The React~18 application wraps all routes in an \texttt{AuthProvider} context for global authentication state, a \texttt{ResourceProvider} for resource library data, and a \texttt{PageErrorBoundary} for graceful error recovery. Role-based route protection ensures providers and clients are directed to separate dashboards:

\begin{lstlisting}[caption={Role-based route protection (\texttt{App.tsx})}]
<ProtectedRoute allowedRoles={['provider']}>
  <ErrorBoundary>
    <ProviderDashboard />
  </ErrorBoundary>
</ProtectedRoute>

<ProtectedRoute allowedRoles={['client']}>
  <ErrorBoundary>
    <ClientDashboard />
  </ErrorBoundary>
</ProtectedRoute>
\end{lstlisting}

\subsubsection{Client Check-in Component --- \texttt{src/components/client/ClientCheckIns.tsx}}

The check-in form captures mood via a slider (1--10), physical symptoms via a checkbox list of ten predefined options, optional notes, and a provider-sharing toggle:

\begin{lstlisting}[caption={Symptom definitions (\texttt{ClientCheckIns.tsx})}]
const SYMPTOMS: Array<{
  id: PhysicalSymptom; label: string
}> = [
  { id: 'fatigue', label: 'Fatigue / low energy' },
  { id: 'sleep_issues', label: 'Sleep challenges' },
  { id: 'appetite_changes', label: 'Appetite changes' },
  { id: 'anxiety', label: 'Anxiety / worry' },
  { id: 'pain', label: 'Pain' },
  { id: 'headache', label: 'Headache' },
  { id: 'nausea', label: 'Nausea' },
  { id: 'dizziness', label: 'Dizziness' },
  { id: 'breast_soreness',
    label: 'Breast / chest soreness' },
  { id: 'bleeding', label: 'Bleeding' },
];
\end{lstlisting}

\subsubsection{Intake Wizard --- \texttt{src/components/intake/ClientIntakeWizard.tsx}}

The five-step wizard validates each section with Zod schemas before allowing progression:

\begin{lstlisting}[caption={Intake wizard Zod validation schemas (\texttt{ClientIntakeWizard.tsx})}]
type StepId = 'personal' | 'birth' | 'feeding'
            | 'support' | 'health';

const schemas: Record<StepId, z.ZodTypeAny> = {
  personal: z.object({
    partnerName: optionalShort(120),
    address: z.object({
      street, city, state, zipCode
    }).optional(),
    emergencyContact: z.object({
      name, phone, relationship
    }).optional(),
    communicationPreferences: z.object({
      preferredContactMethod, bestTimeToContact
    }).optional(),
  }),
  birth: z.object({
    isFirstBaby: z.boolean().optional(),
    birthExperience: z.object({
      birthType, birthLocation, laborDuration
    }).optional(),
  }),
  // feeding, support, health schemas follow same pattern
};
\end{lstlisting}

\newpage

\subsection{Source Code File Inventory}

\begin{table}[H]
\centering
\small
\begin{tabularx}{\textwidth}{|l|l|c|>{\raggedright\arraybackslash}X|}
\hline
\rowcolor{TableHeader}\textcolor{TableHeaderText}{\textbf{Layer}} & \textcolor{TableHeaderText}{\textbf{Directory}} & \textcolor{TableHeaderText}{\textbf{Files}} & \textcolor{TableHeaderText}{\textbf{Description}} \\
\hline
Backend Routes & \texttt{backend/src/routes/} & 20 & API endpoint handlers for all domains \\
\hline
\rowcolor{TableRowAlt}
Backend Models & \texttt{backend/src/models/} & 19 & Mongoose schemas and interfaces \\
\hline
Backend Services & \texttt{backend/src/services/} & 9 & GridFS, email, push, trends, cache, appointment notifications, reminders, rate limiter, socket manager \\
\hline
\rowcolor{TableRowAlt}
Backend Config & \texttt{backend/src/config/} & 1 & Passport authentication strategies \\
\hline
Backend Utils & \texttt{backend/src/utils/} & 4 & Token utils, logger, query optimization, error helpers \\
\hline
\rowcolor{TableRowAlt}
Frontend Pages & \texttt{Lunara/src/pages/} & 10 & Route-level entry points \\
\hline
Frontend Components & \texttt{Lunara/src/components/} & 47 & Provider, client, documents, intake, blog, resource, UI \\
\hline
\rowcolor{TableRowAlt}
Frontend Contexts & \texttt{Lunara/src/contexts/} & 4 & AuthContext, ResourceContext, useAuth, useResource \\
\hline
Frontend Services & \texttt{Lunara/src/services/} & 13 & API service layer for each domain \\
\hline
\end{tabularx}
\caption{Source code file inventory by architectural layer}
\end{table}

\newpage

% ============================================================
% SECTION 5: SCOPE ADJUSTMENT
% ============================================================
\section{Milestone 4 Special Note --- Scope Adjustment Context}

As documented in sprint planning, four planned features were formally deferred to a post-launch roadmap after discussion with the mentor and instructor: FR12 (Digital Journaling), FR13 (Daily Insights / Horoscopes), FR14 (Sleep \& Feeding Trackers), and FR15 (AI-Powered Note Summarization). The justification for each deferral was recorded in the sprint documentation. These were lower-priority features from the original requirements that fell outside the timeline after the core functional requirements (FR1--FR11) consumed more development effort than initially estimated due to the depth of integration required across the full stack. Where scope changed, execution focused on preserving the core functional requirements and delivery quality. This milestone established the integrated coding baseline that enabled formal Milestone~5 testing and Milestone~6 release packaging.

% ============================================================
% SECTION 6: DATA DICTIONARY
% ============================================================
\section{Data Dictionary}

The data dictionary defines every persistent entity in the LUNARA MongoDB database. Each table lists the field name, data type, constraints (required, unique, default, min/max), and references to other collections. All models use Mongoose and reside in \texttt{backend/src/models/}.

% -- helper: header row used in firsthead and head --
% NOTE: \ddheaderrow cannot contain & chars inside \newcommand for longtable,
% so we define a helper that writes the full header via \multicolumn trick.
% Instead we just inline the header rows in each table.

% ---------- User ----------
\needspace{10\baselineskip}
\subsection{User}

{\small
\begin{longtable}{|>{\ttfamily\raggedright\arraybackslash}p{3.8cm}|>{\raggedright\arraybackslash}p{2.2cm}|>{\raggedright\arraybackslash}p{2cm}|>{\raggedright\arraybackslash}p{4.5cm}|}
\caption{User --- \texttt{users} collection (\texttt{models/User.ts})}\label{tab:dd-user}\\
\hline
\rowcolor{TableHeader}\textcolor{TableHeaderText}{\textbf{Field}} & \textcolor{TableHeaderText}{\textbf{Type}} & \textcolor{TableHeaderText}{\textbf{Required}} & \textcolor{TableHeaderText}{\textbf{Notes}} \\
\hline
\endfirsthead
\multicolumn{4}{l}{\small\itshape\color{gray} User (continued)}\\[0.5em]
\hline
\rowcolor{TableHeader}\textcolor{TableHeaderText}{\textbf{Field}} & \textcolor{TableHeaderText}{\textbf{Type}} & \textcolor{TableHeaderText}{\textbf{Required}} & \textcolor{TableHeaderText}{\textbf{Notes}} \\
\hline
\endhead
\hline \multicolumn{4}{r}{\small\itshape\color{gray} Continued on next page}\\ \endfoot
\hline \endlastfoot
firstName & String & Yes & Max 50 characters \\ \hline
lastName & String & Yes & Max 50 characters \\ \hline
email & String & Yes & Unique, lowercase, validated \\ \hline
password & String & No & Min 8 chars, bcrypt-hashed \\ \hline
role & Enum & Yes & client $\vert$ provider $\vert$ admin; default \texttt{client} \\ \hline
isEmailVerified & Boolean & No & Default \texttt{false} \\ \hline
oauthProviders & Array & No & Google OAuth data \\ \hline
profile.phone & String & No & Phone validation \\ \hline
profile.timezone & String & No & Default \texttt{America/Phoenix} \\ \hline
mfaEnabled & Boolean & No & Default \texttt{false} \\ \hline
mfaSecret & String & No & TOTP secret \\ \hline
mfaBackupCodes & [String] & No & One-time recovery codes \\ \hline
refreshTokens & [Object] & No & Token + createdAt; 7-day TTL \\ \hline
failedLoginAttempts & Number & No & Default 0 \\ \hline
lockUntil & Date & No & Account lock expiry \\ \hline
intakeCompleted & Boolean & No & Default \texttt{false} \\ \hline
createdAt / updatedAt & Date & Auto & Mongoose timestamps \\ \hline
\end{longtable}
}

% ---------- Client ----------
\needspace{10\baselineskip}
\subsection{Client}

{\small
\begin{longtable}{|>{\ttfamily\raggedright\arraybackslash}p{3.8cm}|>{\raggedright\arraybackslash}p{2.2cm}|>{\raggedright\arraybackslash}p{2cm}|>{\raggedright\arraybackslash}p{4.5cm}|}
\caption{Client --- \texttt{clients} collection (\texttt{models/Client.ts})}\label{tab:dd-client}\\
\hline
\rowcolor{TableHeader}\textcolor{TableHeaderText}{\textbf{Field}} & \textcolor{TableHeaderText}{\textbf{Type}} & \textcolor{TableHeaderText}{\textbf{Required}} & \textcolor{TableHeaderText}{\textbf{Notes}} \\
\hline
\endfirsthead
\multicolumn{4}{l}{\small\itshape\color{gray} Client (continued)}\\[0.5em]
\hline
\rowcolor{TableHeader}\textcolor{TableHeaderText}{\textbf{Field}} & \textcolor{TableHeaderText}{\textbf{Type}} & \textcolor{TableHeaderText}{\textbf{Required}} & \textcolor{TableHeaderText}{\textbf{Notes}} \\
\hline
\endhead
\hline \multicolumn{4}{r}{\small\itshape\color{gray} Continued on next page}\\ \endfoot
\hline \endlastfoot
userId & ObjectId & Yes & Ref: \texttt{User}; unique \\ \hline
birthDate & Date & No & Client date of birth \\ \hline
babyBirthDate & Date & No & Baby's date of birth \\ \hline
dueDate & Date & No & Expected due date \\ \hline
assignedProvider & ObjectId & No & Ref: \texttt{User} \\ \hline
status & Enum & No & active $\vert$ inactive $\vert$ completed \\ \hline
intakeCompleted & Boolean & No & Default \texttt{false} \\ \hline
intakeData & Object & No & Nested: address, emergency contact, birth experience, feeding preferences, support needs, medical history, postpartum goals \\ \hline
postpartumWeek & Number & No & Default 0 \\ \hline
carePlans & [Object] & No & planId (Ref: \texttt{CarePlan}), assignedDate, status \\ \hline
onboardingSteps & Object & No & Booleans: profileCreated, intakeCompleted, providerAssigned, firstContactMade, careplanAssigned \\ \hline
createdAt / updatedAt & Date & Auto & Mongoose timestamps \\ \hline
\end{longtable}
}

% ---------- Provider ----------
\needspace{10\baselineskip}
\subsection{Provider}

{\small
\begin{longtable}{|>{\ttfamily\raggedright\arraybackslash}p{3.8cm}|>{\raggedright\arraybackslash}p{2.2cm}|>{\raggedright\arraybackslash}p{2cm}|>{\raggedright\arraybackslash}p{4.5cm}|}
\caption{Provider --- \texttt{providers} collection (\texttt{models/Provider.ts})}\label{tab:dd-provider}\\
\hline
\rowcolor{TableHeader}\textcolor{TableHeaderText}{\textbf{Field}} & \textcolor{TableHeaderText}{\textbf{Type}} & \textcolor{TableHeaderText}{\textbf{Required}} & \textcolor{TableHeaderText}{\textbf{Notes}} \\
\hline
\endfirsthead
\multicolumn{4}{l}{\small\itshape\color{gray} Provider (continued)}\\[0.5em]
\hline
\rowcolor{TableHeader}\textcolor{TableHeaderText}{\textbf{Field}} & \textcolor{TableHeaderText}{\textbf{Type}} & \textcolor{TableHeaderText}{\textbf{Required}} & \textcolor{TableHeaderText}{\textbf{Notes}} \\
\hline
\endhead
\hline \multicolumn{4}{r}{\small\itshape\color{gray} Continued on next page}\\ \endfoot
\hline \endlastfoot
userId & ObjectId & Yes & Ref: \texttt{User}; unique \\ \hline
professionalInfo & Object & No & certifications (enum array), yearsExperience (0--50), specialties (enum array), bio (max 1000), education, languages \\ \hline
contactInfo & Object & No & businessPhone, businessEmail, website, socialMedia \\ \hline
serviceAreas & [String] & No & Cities, zip codes, regions \\ \hline
services & [Enum] & No & prenatal\_support, birth\_support, postpartum\_support, overnight\_care, meal\_preparation, etc. \\ \hline
availability & Object & No & isAcceptingClients (default true), maxClients (1--20, default 8), weekly schedule, holidays, vacation dates \\ \hline
clients & [Object] & No & clientId (Ref: \texttt{User}), assignedDate, status, serviceType \\ \hline
pricing & Object & No & hourlyRate, packageRates[], currency (default USD) \\ \hline
status & Enum & No & active $\vert$ inactive $\vert$ pending\_approval $\vert$ suspended \\ \hline
metrics & Object & No & totalClients, averageRating (0--5), totalReviews, responseTime (hrs) \\ \hline
createdAt / updatedAt & Date & Auto & Mongoose timestamps \\ \hline
\end{longtable}
}

% ---------- Appointment ----------
\needspace{10\baselineskip}
\subsection{Appointment}

{\small
\begin{longtable}{|>{\ttfamily\raggedright\arraybackslash}p{3.8cm}|>{\raggedright\arraybackslash}p{2.2cm}|>{\raggedright\arraybackslash}p{2cm}|>{\raggedright\arraybackslash}p{4.5cm}|}
\caption{Appointment --- \texttt{appointments} collection (\texttt{models/Appointment.ts})}\label{tab:dd-appointment}\\
\hline
\rowcolor{TableHeader}\textcolor{TableHeaderText}{\textbf{Field}} & \textcolor{TableHeaderText}{\textbf{Type}} & \textcolor{TableHeaderText}{\textbf{Required}} & \textcolor{TableHeaderText}{\textbf{Notes}} \\
\hline
\endfirsthead
\multicolumn{4}{l}{\small\itshape\color{gray} Appointment (continued)}\\[0.5em]
\hline
\rowcolor{TableHeader}\textcolor{TableHeaderText}{\textbf{Field}} & \textcolor{TableHeaderText}{\textbf{Type}} & \textcolor{TableHeaderText}{\textbf{Required}} & \textcolor{TableHeaderText}{\textbf{Notes}} \\
\hline
\endhead
\hline \multicolumn{4}{r}{\small\itshape\color{gray} Continued on next page}\\ \endfoot
\hline \endlastfoot
clientId & ObjectId & Yes & Ref: \texttt{User} \\ \hline
providerId & ObjectId & Yes & Ref: \texttt{User} \\ \hline
startTime & Date & Yes & --- \\ \hline
endTime & Date & Yes & --- \\ \hline
status & Enum & No & requested $\vert$ confirmed $\vert$ scheduled $\vert$ completed $\vert$ cancelled \\ \hline
type & Enum & No & virtual $\vert$ in\_person; default \texttt{virtual} \\ \hline
notes & String & No & Free-text notes \\ \hline
requestedBy & ObjectId & No & Ref: \texttt{User} \\ \hline
cancellationReason & String & No & Max 500 characters \\ \hline
cancelledBy & ObjectId & No & Ref: \texttt{User} \\ \hline
availabilitySlotId & ObjectId & No & Ref: \texttt{AvailabilitySlot} \\ \hline
createdAt / updatedAt & Date & Auto & Mongoose timestamps \\ \hline
\end{longtable}
}

% ---------- AvailabilitySlot ----------
\needspace{10\baselineskip}
\subsection{AvailabilitySlot}

{\small
\begin{longtable}{|>{\ttfamily\raggedright\arraybackslash}p{3.8cm}|>{\raggedright\arraybackslash}p{2.2cm}|>{\raggedright\arraybackslash}p{2cm}|>{\raggedright\arraybackslash}p{4.5cm}|}
\caption{AvailabilitySlot --- \texttt{availabilityslots} collection (\texttt{models/AvailabilitySlot.ts})}\label{tab:dd-avail}\\
\hline
\rowcolor{TableHeader}\textcolor{TableHeaderText}{\textbf{Field}} & \textcolor{TableHeaderText}{\textbf{Type}} & \textcolor{TableHeaderText}{\textbf{Required}} & \textcolor{TableHeaderText}{\textbf{Notes}} \\
\hline
\endfirsthead
\multicolumn{4}{l}{\small\itshape\color{gray} AvailabilitySlot (continued)}\\[0.5em]
\hline
\rowcolor{TableHeader}\textcolor{TableHeaderText}{\textbf{Field}} & \textcolor{TableHeaderText}{\textbf{Type}} & \textcolor{TableHeaderText}{\textbf{Required}} & \textcolor{TableHeaderText}{\textbf{Notes}} \\
\hline
\endhead
\hline \multicolumn{4}{r}{\small\itshape\color{gray} Continued on next page}\\ \endfoot
\hline \endlastfoot
providerId & ObjectId & Yes & Ref: \texttt{User} \\ \hline
date & Date & Yes & Slot date \\ \hline
startTime & String & Yes & HH:mm format \\ \hline
endTime & String & Yes & HH:mm format \\ \hline
isBooked & Boolean & No & Default \texttt{false} \\ \hline
appointmentId & ObjectId & No & Ref: \texttt{Appointment} \\ \hline
recurring & Boolean & No & Default \texttt{false} \\ \hline
dayOfWeek & Number & No & 0--6 (Sun--Sat) \\ \hline
createdAt / updatedAt & Date & Auto & Mongoose timestamps \\ \hline
\end{longtable}
}

% ---------- CarePlan ----------
\needspace{10\baselineskip}
\subsection{CarePlan}

{\small
\begin{longtable}{|>{\ttfamily\raggedright\arraybackslash}p{3.8cm}|>{\raggedright\arraybackslash}p{2.2cm}|>{\raggedright\arraybackslash}p{2cm}|>{\raggedright\arraybackslash}p{4.5cm}|}
\caption{CarePlan --- \texttt{careplans} collection (\texttt{models/CarePlan.ts})}\label{tab:dd-careplan}\\
\hline
\rowcolor{TableHeader}\textcolor{TableHeaderText}{\textbf{Field}} & \textcolor{TableHeaderText}{\textbf{Type}} & \textcolor{TableHeaderText}{\textbf{Required}} & \textcolor{TableHeaderText}{\textbf{Notes}} \\
\hline
\endfirsthead
\multicolumn{4}{l}{\small\itshape\color{gray} CarePlan (continued)}\\[0.5em]
\hline
\rowcolor{TableHeader}\textcolor{TableHeaderText}{\textbf{Field}} & \textcolor{TableHeaderText}{\textbf{Type}} & \textcolor{TableHeaderText}{\textbf{Required}} & \textcolor{TableHeaderText}{\textbf{Notes}} \\
\hline
\endhead
\hline \multicolumn{4}{r}{\small\itshape\color{gray} Continued on next page}\\ \endfoot
\hline \endlastfoot
clientId & ObjectId & Yes & Ref: \texttt{User} \\ \hline
providerId & ObjectId & Yes & Ref: \texttt{User} \\ \hline
templateId & ObjectId & No & Ref: \texttt{CarePlanTemplate} \\ \hline
title & String & Yes & Max 200 characters \\ \hline
description & String & No & Max 2000 characters \\ \hline
sections & [Object] & Yes & Each section: title, description, milestones[] \\ \hline
status & Enum & No & active $\vert$ completed $\vert$ paused $\vert$ archived \\ \hline
startDate & Date & No & Default \texttt{Date.now} \\ \hline
endDate & Date & No & --- \\ \hline
progress & Number & No & 0--100 (percentage) \\ \hline
createdAt / updatedAt & Date & Auto & Mongoose timestamps \\ \hline
\end{longtable}
}

% ---------- CheckIn ----------
\needspace{10\baselineskip}
\subsection{CheckIn}

{\small
\begin{longtable}{|>{\ttfamily\raggedright\arraybackslash}p{3.8cm}|>{\raggedright\arraybackslash}p{2.2cm}|>{\raggedright\arraybackslash}p{2cm}|>{\raggedright\arraybackslash}p{4.5cm}|}
\caption{CheckIn --- \texttt{checkins} collection (\texttt{models/CheckIn.ts})}\label{tab:dd-checkin}\\
\hline
\rowcolor{TableHeader}\textcolor{TableHeaderText}{\textbf{Field}} & \textcolor{TableHeaderText}{\textbf{Type}} & \textcolor{TableHeaderText}{\textbf{Required}} & \textcolor{TableHeaderText}{\textbf{Notes}} \\
\hline
\endfirsthead
\multicolumn{4}{l}{\small\itshape\color{gray} CheckIn (continued)}\\[0.5em]
\hline
\rowcolor{TableHeader}\textcolor{TableHeaderText}{\textbf{Field}} & \textcolor{TableHeaderText}{\textbf{Type}} & \textcolor{TableHeaderText}{\textbf{Required}} & \textcolor{TableHeaderText}{\textbf{Notes}} \\
\hline
\endhead
\hline \multicolumn{4}{r}{\small\itshape\color{gray} Continued on next page}\\ \endfoot
\hline \endlastfoot
userId & ObjectId & Yes & Ref: \texttt{User} \\ \hline
date & Date & Yes & Check-in date \\ \hline
moodScore & Number & Yes & 1--10 scale \\ \hline
physicalSymptoms & [Enum] & No & fatigue, sleep\_issues, appetite\_changes, anxiety, pain, headache, nausea, dizziness, breast\_soreness, bleeding \\ \hline
notes & String & No & Max 2000 characters \\ \hline
sharedWithProvider & Boolean & No & Default \texttt{false} \\ \hline
providerReviewed & Boolean & No & Default \texttt{false} \\ \hline
createdAt / updatedAt & Date & Auto & Mongoose timestamps \\ \hline
\end{longtable}
}

% ---------- Message ----------
\needspace{10\baselineskip}
\subsection{Message}

{\small
\begin{longtable}{|>{\ttfamily\raggedright\arraybackslash}p{3.8cm}|>{\raggedright\arraybackslash}p{2.2cm}|>{\raggedright\arraybackslash}p{2cm}|>{\raggedright\arraybackslash}p{4.5cm}|}
\caption{Message --- \texttt{messages} collection (\texttt{models/Message.ts})}\label{tab:dd-message}\\
\hline
\rowcolor{TableHeader}\textcolor{TableHeaderText}{\textbf{Field}} & \textcolor{TableHeaderText}{\textbf{Type}} & \textcolor{TableHeaderText}{\textbf{Required}} & \textcolor{TableHeaderText}{\textbf{Notes}} \\
\hline
\endfirsthead
\multicolumn{4}{l}{\small\itshape\color{gray} Message (continued)}\\[0.5em]
\hline
\rowcolor{TableHeader}\textcolor{TableHeaderText}{\textbf{Field}} & \textcolor{TableHeaderText}{\textbf{Type}} & \textcolor{TableHeaderText}{\textbf{Required}} & \textcolor{TableHeaderText}{\textbf{Notes}} \\
\hline
\endhead
\hline \multicolumn{4}{r}{\small\itshape\color{gray} Continued on next page}\\ \endfoot
\hline \endlastfoot
conversationId & ObjectId & Yes & Groups messages in a thread \\ \hline
sender & ObjectId & Yes & Ref: \texttt{User} \\ \hline
receiver & ObjectId & Yes & Ref: \texttt{User} \\ \hline
content & String & Yes & Max 10\,000 characters \\ \hline
title & String & No & Optional subject line \\ \hline
read & Boolean & No & Default \texttt{false} \\ \hline
type & Enum & No & text $\vert$ image $\vert$ file $\vert$ system \\ \hline
createdAt & Date & Auto & Mongoose timestamp \\ \hline
\end{longtable}
}

% ---------- ClientDocument ----------
\needspace{10\baselineskip}
\subsection{ClientDocument}

{\small
\begin{longtable}{|>{\ttfamily\raggedright\arraybackslash}p{3.8cm}|>{\raggedright\arraybackslash}p{2.2cm}|>{\raggedright\arraybackslash}p{2cm}|>{\raggedright\arraybackslash}p{4.5cm}|}
\caption{ClientDocument --- \texttt{clientdocuments} collection (\texttt{models/ClientDocument.ts})}\label{tab:dd-clientdoc}\\
\hline
\rowcolor{TableHeader}\textcolor{TableHeaderText}{\textbf{Field}} & \textcolor{TableHeaderText}{\textbf{Type}} & \textcolor{TableHeaderText}{\textbf{Required}} & \textcolor{TableHeaderText}{\textbf{Notes}} \\
\hline
\endfirsthead
\multicolumn{4}{l}{\small\itshape\color{gray} ClientDocument (continued)}\\[0.5em]
\hline
\rowcolor{TableHeader}\textcolor{TableHeaderText}{\textbf{Field}} & \textcolor{TableHeaderText}{\textbf{Type}} & \textcolor{TableHeaderText}{\textbf{Required}} & \textcolor{TableHeaderText}{\textbf{Notes}} \\
\hline
\endhead
\hline \multicolumn{4}{r}{\small\itshape\color{gray} Continued on next page}\\ \endfoot
\hline \endlastfoot
title & String & Yes & Max 200 characters \\ \hline
documentType & Enum & Yes & emotional-survey, health-assessment, personal-assessment, feeding-log, sleep-log, mood-check-in, recovery-notes, progress-photo, other \\ \hline
uploadedBy & ObjectId & Yes & Ref: \texttt{User} \\ \hline
assignedProvider & ObjectId & No & Ref: \texttt{User} \\ \hline
files & [Object] & No & cloudinaryUrl, originalFileName, fileType, fileSize, uploadDate \\ \hline
submissionStatus & Enum & No & draft $\vert$ submitted-to-provider $\vert$ reviewed-by-provider $\vert$ completed \\ \hline
formData & Mixed & No & Survey responses, assessment scores \\ \hline
privacyLevel & Enum & No & client-only $\vert$ client-and-provider $\vert$ care-team \\ \hline
dueDate & Date & No & --- \\ \hline
createdAt / updatedAt & Date & Auto & Mongoose timestamps \\ \hline
\end{longtable}
}

% ---------- BlogPost ----------
\needspace{10\baselineskip}
\subsection{BlogPost}

{\small
\begin{longtable}{|>{\ttfamily\raggedright\arraybackslash}p{3.8cm}|>{\raggedright\arraybackslash}p{2.2cm}|>{\raggedright\arraybackslash}p{2cm}|>{\raggedright\arraybackslash}p{4.5cm}|}
\caption{BlogPost --- \texttt{blogposts} collection (\texttt{models/BlogPost.ts})}\label{tab:dd-blog}\\
\hline
\rowcolor{TableHeader}\textcolor{TableHeaderText}{\textbf{Field}} & \textcolor{TableHeaderText}{\textbf{Type}} & \textcolor{TableHeaderText}{\textbf{Required}} & \textcolor{TableHeaderText}{\textbf{Notes}} \\
\hline
\endfirsthead
\multicolumn{4}{l}{\small\itshape\color{gray} BlogPost (continued)}\\[0.5em]
\hline
\rowcolor{TableHeader}\textcolor{TableHeaderText}{\textbf{Field}} & \textcolor{TableHeaderText}{\textbf{Type}} & \textcolor{TableHeaderText}{\textbf{Required}} & \textcolor{TableHeaderText}{\textbf{Notes}} \\
\hline
\endhead
\hline \multicolumn{4}{r}{\small\itshape\color{gray} Continued on next page}\\ \endfoot
\hline \endlastfoot
title & String & Yes & Max 200 characters \\ \hline
slug & String & Yes & Unique, lowercase \\ \hline
excerpt & String & Yes & Max 500 characters \\ \hline
content & String & Yes & Max 500\,KB \\ \hline
author & ObjectId & Yes & Ref: \texttt{User} \\ \hline
featuredImage & String & No & Cloudinary URL \\ \hline
tags & [String] & No & Indexed for search \\ \hline
category & String & Yes & Indexed \\ \hline
isPublished & Boolean & No & Default \texttt{false} \\ \hline
publishDate & Date & No & --- \\ \hline
readTime & Number & No & 1--60 minutes \\ \hline
viewCount & Number & No & Default 0 \\ \hline
createdAt / updatedAt & Date & Auto & Mongoose timestamps \\ \hline
\end{longtable}
}

% ---------- Resource ----------
\needspace{10\baselineskip}
\subsection{Resource}

{\small
\begin{longtable}{|>{\ttfamily\raggedright\arraybackslash}p{3.8cm}|>{\raggedright\arraybackslash}p{2.2cm}|>{\raggedright\arraybackslash}p{2cm}|>{\raggedright\arraybackslash}p{4.5cm}|}
\caption{Resource --- \texttt{resources} collection (\texttt{models/Resources.ts})}\label{tab:dd-resource}\\
\hline
\rowcolor{TableHeader}\textcolor{TableHeaderText}{\textbf{Field}} & \textcolor{TableHeaderText}{\textbf{Type}} & \textcolor{TableHeaderText}{\textbf{Required}} & \textcolor{TableHeaderText}{\textbf{Notes}} \\
\hline
\endfirsthead
\multicolumn{4}{l}{\small\itshape\color{gray} Resource (continued)}\\[0.5em]
\hline
\rowcolor{TableHeader}\textcolor{TableHeaderText}{\textbf{Field}} & \textcolor{TableHeaderText}{\textbf{Type}} & \textcolor{TableHeaderText}{\textbf{Required}} & \textcolor{TableHeaderText}{\textbf{Notes}} \\
\hline
\endhead
\hline \multicolumn{4}{r}{\small\itshape\color{gray} Continued on next page}\\ \endfoot
\hline \endlastfoot
title & String & Yes & Max 200 characters \\ \hline
description & String & Yes & Max 1000 characters \\ \hline
content & String & Yes & Max 500\,KB \\ \hline
category & String & Yes & Indexed \\ \hline
tags & [String] & No & Max 50 tags \\ \hline
author & ObjectId & Yes & Ref: \texttt{User} \\ \hline
targetWeeks & [Number] & No & Postpartum weeks 1--52 \\ \hline
targetPregnancyWeeks & [Number] & No & Pregnancy weeks 1--42 \\ \hline
difficulty & Enum & No & beginner $\vert$ intermediate $\vert$ advanced \\ \hline
isPublished & Boolean & No & Default \texttt{false} \\ \hline
createdAt / updatedAt & Date & Auto & Mongoose timestamps \\ \hline
\end{longtable}
}

% ---------- Remaining models ----------
\needspace{10\baselineskip}
\subsection{Supporting Models}

The following additional models support version history, public inquiries, push notifications, and user interactions:

\begin{itemize}[nosep]
  \item \textbf{CarePlanTemplate} --- Reusable care plan blueprints with sections and milestone definitions. Fields: name, description, targetCondition, sections[], isActive, createdBy (Ref: User).
  \item \textbf{BlogPostVersion} --- Immutable audit trail for blog edits. Fields: blogPostId (Ref: BlogPost), versionNumber, full content snapshot, changedBy (Ref: User), changeReason.
  \item \textbf{ClientDocumentVersion} --- Version history for client documents. Fields: documentId (Ref: ClientDocument), versionNumber, full document snapshot, changedBy (Ref: User).
  \item \textbf{ResourceVersion} --- Version history for resources. Fields: resourceId (Ref: Resource), versionNumber, full content snapshot, changedBy (Ref: User).
  \item \textbf{Category} --- Hierarchical taxonomy with self-referential parent/subcategory structure. Fields: name, description, parentCategory (Ref: Category), subcategories[].
  \item \textbf{Inquiry} --- Public contact-form submissions. Fields: name, email, phone, message, dueDate, status (new $\vert$ contacted $\vert$ converted $\vert$ closed), respondedBy (Ref: User).
  \item \textbf{PushSubscription} --- Web Push API subscription records. Fields: userId (Ref: User), endpoint (unique), keys (p256dh, auth).
  \item \textbf{UserResourceInteraction} --- Tracks views, downloads, favorites, and ratings. Fields: user (Ref: User), resource (Ref: Resource), interactionType, rating.
\end{itemize}

\needspace{10\baselineskip}
\subsection{Entity Relationship Summary}

\textbf{19 collections total.} The \texttt{User} model serves as the central identity entity referenced by nearly every other collection. Key relationships:

\begin{itemize}[nosep]
  \item \texttt{User} $\longleftrightarrow$ \texttt{Client} / \texttt{Provider} (1:1 via \texttt{userId})
  \item \texttt{Client} $\longrightarrow$ \texttt{Provider} (M:1 via \texttt{assignedProvider})
  \item \texttt{Provider.clients[]} $\longrightarrow$ \texttt{User} (embedded reference array)
  \item \texttt{Appointment} $\longrightarrow$ \texttt{User} $\times$ 2 (clientId, providerId)
  \item \texttt{Appointment} $\longrightarrow$ \texttt{AvailabilitySlot} (optional link)
  \item \texttt{CarePlan} $\longrightarrow$ \texttt{User} $\times$ 2 + optional \texttt{CarePlanTemplate}
  \item \texttt{Message} $\longrightarrow$ \texttt{User} $\times$ 2 (sender, receiver)
  \item Version models $\longrightarrow$ parent document + \texttt{User} (changedBy)
\end{itemize}

\newpage

% ============================================================
% SECTION 7: APPLICATION DEMONSTRATION
% ============================================================
\section{Application Demonstration}

The Milestone~4 application demonstration is a screencast (5--10 minutes per GCU policies) that demonstrates the functionality working up to this point. The demonstration covers the public landing page, authentication flows (register, login, MFA), the provider dashboard with client management and messaging, the client dashboard with intake wizard and check-ins, appointment scheduling, document upload and review, and the resource library.

\subsection{Video Demonstration}

The full video demonstration for Milestone~4 is available at the following link:

\begin{center}
\url{https://www.youtube.com/watch?v=EY6ncwOMFUQ}
\end{center}

\newpage

% ============================================================
% SECTION 7: MODULE TEST CASES
% ============================================================
\section{Module Test Cases (Test Plan)}

The following test cases were developed and executed during the Milestone~4 coding iteration. Each test case maps to one or more functional requirements from the Traceability Matrix. Testing was performed manually against the deployed staging environment (\texttt{lunara.onrender.com}) and locally via \texttt{npm run dev}. Automated unit tests supplement manual verification and are maintained in the repository under \texttt{backend/tests/} and \texttt{Lunara/src/tests/}.

\subsection{Test Case Summary}

{\small
\begin{longtable}{|c|c|>{\raggedright\arraybackslash}p{3cm}|>{\raggedright\arraybackslash}p{6.5cm}|}
\caption{Module Test Case Summary (TC-001 through TC-020)}\label{tab:test-summary}\\
\hline
\rowcolor{TableHeader}\textcolor{TableHeaderText}{\textbf{ID}} & \textcolor{TableHeaderText}{\textbf{Priority}} & \textcolor{TableHeaderText}{\textbf{Module}} & \textcolor{TableHeaderText}{\textbf{Test Objective}} \\
\hline
\endfirsthead
\multicolumn{4}{l}{\small\itshape\color{gray} Table~\ref{tab:test-summary} continued from previous page}\\[0.5em]
\hline
\rowcolor{TableHeader}\textcolor{TableHeaderText}{\textbf{ID}} & \textcolor{TableHeaderText}{\textbf{Priority}} & \textcolor{TableHeaderText}{\textbf{Module}} & \textcolor{TableHeaderText}{\textbf{Test Objective}} \\
\hline
\endhead
\hline
\multicolumn{4}{r}{\small\itshape\color{gray} Continued on next page}\\
\endfoot
\hline
\endlastfoot

TC-001 & P1 & Landing Page & Verify the public landing page loads with hero, offerings accordion, inquiry form, and footer \\
\hline
\rowcolor{TableRowAlt}
TC-002 & P1 & Contact Form & Verify the inquiry form validates required fields and submits to \texttt{POST /api/public/contact} \\
\hline
TC-003 & P1 & User Registration & Verify a new user can register with valid email, password (8+ chars), first/last name, and role \\
\hline
\rowcolor{TableRowAlt}
TC-004 & P1 & Token Refresh & Verify an expired access token is replaced via \texttt{POST /api/auth/refresh-token} \\
\hline
TC-005 & P1 & Protected Route & Verify requests without a valid JWT return \texttt{401 Unauthorized} \\
\hline
\rowcolor{TableRowAlt}
TC-006 & P2 & MFA Setup & Verify TOTP setup returns a QR code and that a valid 6-digit code completes enrollment \\
\hline
TC-007 & P1 & Provider Dashboard & Verify the provider dashboard renders overview metrics, quick actions, upcoming appointments \\
\hline
\rowcolor{TableRowAlt}
TC-008 & P1 & Client Dashboard & Verify the client dashboard renders mood orb, document count, resource count \\
\hline
TC-009 & P1 & Intake Wizard & Verify the five-step wizard validates each section (Zod) and submits the complete intake \\
\hline
\rowcolor{TableRowAlt}
TC-010 & P1 & Messaging & Verify real-time messaging via Socket.IO between client and provider \\
\hline
TC-011 & P1 & Appointment CRUD & Verify appointments can be created, approved, confirmed, and cancelled \\
\hline
\rowcolor{TableRowAlt}
TC-012 & P1 & Resource Library & Verify resources are browsable, filterable, and display full content \\
\hline
TC-013 & P1 & Check-in Submit & Verify a daily check-in with mood, symptoms, and notes saves and returns alerts \\
\hline
\rowcolor{TableRowAlt}
TC-014 & P2 & Trend Compute & Verify the trend service computes mood averages, direction, and streak alerts \\
\hline
TC-015 & P1 & Client Management & Verify a provider can create, assign, and unassign clients \\
\hline
\rowcolor{TableRowAlt}
TC-016 & P2 & Care Plan CRUD & Verify care plans can be created, milestones completed, and progress auto-computes \\
\hline
TC-017 & P1 & Blog Publish & Verify blog posts can be created and published posts appear publicly \\
\hline
\rowcolor{TableRowAlt}
TC-018 & P1 & Password Hashing & Verify passwords are stored as bcrypt hashes and never returned in plaintext \\
\hline
TC-019 & P1 & Validation Reject & Verify malformed input is rejected with \texttt{422} \\
\hline
\rowcolor{TableRowAlt}
TC-020 & P1 & Socket Auth & Verify Socket.IO connections without a valid JWT are rejected \\
\hline
\end{longtable}
}

\newpage

\subsection{Detailed Test Case Specifications}

\needspace{8\baselineskip}
\subsubsection*{TC-001: Landing Page Load}

\begin{table}[H]
\centering
\small
\begin{tabularx}{\textwidth}{|c|>{\raggedright\arraybackslash}X|>{\raggedright\arraybackslash}p{2cm}|>{\raggedright\arraybackslash}X|c|}
\hline
\rowcolor{TableHeader}\textcolor{TableHeaderText}{\textbf{Step}} & \textcolor{TableHeaderText}{\textbf{Action}} & \textcolor{TableHeaderText}{\textbf{Test Data}} & \textcolor{TableHeaderText}{\textbf{Expected Result}} & \textcolor{TableHeaderText}{\textbf{P/F}} \\
\hline
1 & Navigate to \texttt{https://www.lunaracare.org} & --- & Page loads with hero section, navigation, and footer & Pass \\
\hline
\rowcolor{TableRowAlt}
2 & Scroll to offerings section & --- & Accordion items expand/collapse on click & Pass \\
\hline
3 & Scroll to inquiry form & --- & Form displays First Name, Last Name, Email, Message fields & Pass \\
\hline
\rowcolor{TableRowAlt}
4 & Resize browser to 375px width & --- & Layout adapts to mobile; no horizontal scroll & Pass \\
\hline
\end{tabularx}
\end{table}

\needspace{8\baselineskip}
\subsubsection*{TC-003: User Registration Flow}

\begin{table}[H]
\centering
\small
\begin{tabularx}{\textwidth}{|c|>{\raggedright\arraybackslash}X|>{\raggedright\arraybackslash}p{3cm}|>{\raggedright\arraybackslash}X|c|}
\hline
\rowcolor{TableHeader}\textcolor{TableHeaderText}{\textbf{Step}} & \textcolor{TableHeaderText}{\textbf{Action}} & \textcolor{TableHeaderText}{\textbf{Test Data}} & \textcolor{TableHeaderText}{\textbf{Expected Result}} & \textcolor{TableHeaderText}{\textbf{P/F}} \\
\hline
1 & Navigate to \texttt{/register} & --- & Registration form renders with all fields & Pass \\
\hline
\rowcolor{TableRowAlt}
2 & Submit with missing email & \texttt{firstName: "Test", lastName: "User", password: "Test1234!"} & Validation error displayed; form not submitted & Pass \\
\hline
3 & Submit with password < 8 chars & \texttt{email: "test@ex.com", password: "short"} & Validation error: ``Password must be at least 8 characters'' & Pass \\
\hline
\rowcolor{TableRowAlt}
4 & Submit with valid data & \texttt{firstName: "New", lastName: "Client", role: "client"} & Account created; verification email sent; redirect to login & Pass \\
\hline
\end{tabularx}
\end{table}

\needspace{8\baselineskip}
\subsubsection*{TC-010: Message Send/Receive}

\begin{table}[H]
\centering
\small
\begin{tabularx}{\textwidth}{|c|>{\raggedright\arraybackslash}X|>{\raggedright\arraybackslash}p{3cm}|>{\raggedright\arraybackslash}X|c|}
\hline
\rowcolor{TableHeader}\textcolor{TableHeaderText}{\textbf{Step}} & \textcolor{TableHeaderText}{\textbf{Action}} & \textcolor{TableHeaderText}{\textbf{Test Data}} & \textcolor{TableHeaderText}{\textbf{Expected Result}} & \textcolor{TableHeaderText}{\textbf{P/F}} \\
\hline
1 & Client logs in and navigates to Messages & Valid credentials & Chat interface loads; Socket.IO ``Connected'' indicator & Pass \\
\hline
\rowcolor{TableRowAlt}
2 & Client sends a message & ``Im feeling a bit tired today!'' & Message appears in client chat with timestamp & Pass \\
\hline
3 & Provider opens Message Center & Valid credentials & Provider sees client in sidebar; message in conversation & Pass \\
\hline
\rowcolor{TableRowAlt}
4 & Provider replies & ``hello'' & Reply appears in both views in real time & Pass \\
\hline
\end{tabularx}
\end{table}

\needspace{8\baselineskip}
\subsubsection*{TC-011: Appointment CRUD}

\begin{table}[H]
\centering
\small
\begin{tabularx}{\textwidth}{|c|>{\raggedright\arraybackslash}X|>{\raggedright\arraybackslash}p{3cm}|>{\raggedright\arraybackslash}X|c|}
\hline
\rowcolor{TableHeader}\textcolor{TableHeaderText}{\textbf{Step}} & \textcolor{TableHeaderText}{\textbf{Action}} & \textcolor{TableHeaderText}{\textbf{Test Data}} & \textcolor{TableHeaderText}{\textbf{Expected Result}} & \textcolor{TableHeaderText}{\textbf{P/F}} \\
\hline
1 & Client selects a date on calendar & March 19, 2026 & Date panel shows ``Request time'' option & Pass \\
\hline
\rowcolor{TableRowAlt}
2 & Client fills and submits request & 7:02--9:02 PM, Virtual, ``Casual check in'' & Request submitted; ``Pending approval'' & Pass \\
\hline
3 & Provider views Schedule tab & --- & Pending request with Approve/Decline actions & Pass \\
\hline
\rowcolor{TableRowAlt}
4 & Provider approves the request & --- & Status changes to ``Confirmed''; client sees confirmation & Pass \\
\hline
\end{tabularx}
\end{table}

\needspace{8\baselineskip}
\subsubsection*{TC-013: Check-in Submit}

\begin{table}[H]
\centering
\small
\begin{tabularx}{\textwidth}{|c|>{\raggedright\arraybackslash}X|>{\raggedright\arraybackslash}p{3cm}|>{\raggedright\arraybackslash}X|c|}
\hline
\rowcolor{TableHeader}\textcolor{TableHeaderText}{\textbf{Step}} & \textcolor{TableHeaderText}{\textbf{Action}} & \textcolor{TableHeaderText}{\textbf{Test Data}} & \textcolor{TableHeaderText}{\textbf{Expected Result}} & \textcolor{TableHeaderText}{\textbf{P/F}} \\
\hline
1 & Client selects mood on orb & ``Doing well'' (score 9--10) & Orb turns green; mood label displayed & Pass \\
\hline
\rowcolor{TableRowAlt}
2 & System records check-in & --- & \texttt{POST /api/checkins} returns 201 with data and alerts & Pass \\
\hline
3 & Client selects ``Need support now'' & Score 1--2 & Orb turns red; provider alert generated if streak met & Pass \\
\hline
\end{tabularx}
\end{table}

\needspace{8\baselineskip}
\subsubsection*{TC-017: Blog Publish}

\begin{table}[H]
\centering
\small
\begin{tabularx}{\textwidth}{|c|>{\raggedright\arraybackslash}X|>{\raggedright\arraybackslash}p{3cm}|>{\raggedright\arraybackslash}X|c|}
\hline
\rowcolor{TableHeader}\textcolor{TableHeaderText}{\textbf{Step}} & \textcolor{TableHeaderText}{\textbf{Action}} & \textcolor{TableHeaderText}{\textbf{Test Data}} & \textcolor{TableHeaderText}{\textbf{Expected Result}} & \textcolor{TableHeaderText}{\textbf{P/F}} \\
\hline
1 & Provider navigates to Blog > Create & --- & Blog editor form renders & Pass \\
\hline
\rowcolor{TableRowAlt}
2 & Provider fills in details & Title, category, excerpt, content & All fields accept input; image upload works & Pass \\
\hline
3 & Provider clicks ``Publish'' & --- & Success toast; post in list with ``Published'' status & Pass \\
\hline
\rowcolor{TableRowAlt}
4 & Client navigates to \texttt{/blog} & --- & Published post appears in public blog list & Pass \\
\hline
\end{tabularx}
\end{table}

\newpage

\subsection{Automated Test Execution Results}

In addition to manual test cases, the Lunara project maintains a comprehensive automated test suite covering both the backend API layer and the frontend service/component layer. All tests are executed via Jest and produce coverage reports. The combined suite totals \textbf{38 test suites} and \textbf{374 individual tests}, all passing.

\subsubsection*{Backend Test Results}

The backend test suite covers middleware, services, and utility modules with \textbf{13 test suites} and \textbf{153 tests}. Overall code coverage is \textbf{90.58\% statements}, 75.43\% branches, 87.28\% functions, and 90.56\% lines.

\screenshotfig{Photos/tests/backend.png}{Backend Automated Test Results --- 13 suites, 153 tests passing with 90.58\% statement coverage across middleware, services, and utility modules.}{test-backend}

\subsubsection*{Frontend Test Results}

The frontend test suite covers API clients, services, components, authentication flows, and utility functions with \textbf{25 test suites} and \textbf{221 tests}. Service-layer coverage reaches \textbf{84.27\% statements} with 88.8\% function coverage.

\screenshotfig{Photos/tests/frontend.png}{Frontend Automated Test Results --- 25 suites, 221 tests passing with coverage across API clients, services, components, and authentication flows.}{test-frontend}

\newpage

% ============================================================
% SECTION 8: APPLICATION SCREENSHOTS
% ============================================================
\section{Application Screenshots}

The following screenshots document the working application as demonstrated in the Milestone~4 screencast. They are organized by access level: public pages, provider dashboard workflows, and client dashboard workflows.

% ── Public Website ──
\subsection{Public Website (FR1)}

\screenshotfig{Photos/public/HomePage.png}{Landing Page --- The public homepage features a hero section, the ``Dear Parent'' introduction, client login access, an interactive offerings accordion, and an inquiry contact form.}{landing-page}

\screenshotfig{Photos/public/loginPage.png}{Login Page --- The login page supports email/password authentication with role-based redirection to the appropriate dashboard.}{login-page}

\screenshotfig{Photos/public/publicBlog_page.png}{Public Blog List --- The public-facing blog displays published posts with excerpts, dates, and ``Read More'' links.}{public-blog}

\screenshotfig{Photos/public/publicBlog_detailed.png}{Public Blog Post Detail --- A full blog post view with category tags, author attribution, read time, and view count.}{public-blog-detail}

\newpage

% ── Provider Dashboard ──
\subsection{Provider Dashboard (FR3, FR9, FR11)}

\screenshotfig{Photos/provider/provider_dashboard_overview.png}{Provider Dashboard Overview --- Active client count, pending inquiries, check-ins needing attention, quick action buttons, upcoming appointments, and personal to-do list.}{provider-dashboard}

\screenshotfig{Photos/provider/provider_recent_activity.png}{Provider Recent Activity --- Upcoming appointments, document submissions from clients, and client assignment history.}{provider-activity}

\screenshotfig{Photos/provider/provider_clients_page_names_resolved.png}{Provider Clients Page --- ``My clients'' (assigned) and ``All created clients'' with names, emails, provider assignments, and actions.}{provider-clients}

\screenshotfig{Photos/provider/provider_invite_client_form.png}{Invite New Client Form --- Providers invite new clients by entering first name, last name, and email.}{invite-client}

\screenshotfig{Photos/provider/provider_client_created.png}{New Client Created --- The new client appears in ``All created clients'' with an ``Add to my list'' option.}{client-created}

\newpage

\screenshotfig{Photos/provider/provider_schedule_appointment_complete.png}{Schedule Appointment Modal --- Client and provider selected, start/end times, type (Virtual), and notes.}{schedule-modal}

\screenshotfig{Photos/provider/provider_schedule_calendar.png}{Provider Schedule Calendar --- Monthly calendar view with color-coded appointments and sidebar details.}{provider-calendar}

\screenshotfig{Photos/provider/provider_schedule_availability_slots.png}{Provider Availability Slots --- Providers add availability time slots for specific dates.}{availability-slots}

\newpage

\screenshotfig{Photos/provider/provider_client_document_found.png}{Client Documents Search and Filter --- Search, filter by type, status, and date range with ``Review Now'' action.}{provider-docs}

\screenshotfig{Photos/provider/provider_document_review_panel.png}{Document Review Panel --- Title, author, attached files, internal notes, client-facing feedback, and review status.}{doc-review}

\screenshotfig{Photos/provider/provider_file_view_success.png}{Authenticated File Viewer --- Client-uploaded files viewed through an authenticated blob URL.}{file-viewer}

\screenshotfig{Photos/provider/provider_document_review_completed.png}{Document Review Completed --- Status updates to ``Completed'' and the client is notified.}{review-done}

\newpage

\screenshotfig{Photos/provider/provider_blog_manage_empty.png}{Blog Management --- Manage Blog Posts page with search, filter by category/tag/author, and tabs.}{blog-manage}

\screenshotfig{Photos/provider/provider_blog_create_post.png}{Create New Blog Post --- Blog editor with title, category, excerpt, featured image, tags, and rich text content.}{blog-create}

\screenshotfig{Photos/provider/provider_blog_manage_with_post.png}{Blog Post Published --- Published post in management list with Edit, Publish, and Delete actions.}{blog-published}

\newpage

\screenshotfig{Photos/provider/provider_resource_create_blank.png}{Create New Resource (Blank Form) --- Title, category, description, content, difficulty, week selectors, tags, and publish.}{resource-blank}

\screenshotfig{Photos/provider/provider_resource_create_filled.png}{Create New Resource (Completed Form) --- ``First Trimester: What to Expect'' with weeks 1--12 selected and content filled.}{resource-filled}

\screenshotfig{Photos/provider/provider_resource_in_library.png}{Resource Published in Library --- Resource in library with detail modal showing content, category, difficulty, and week targeting.}{resource-library}

\newpage

\screenshotfig{Photos/provider/provider_profile_edit.png}{Provider Profile Edit --- Personal information, bio, certifications, specialties, languages, services, and availability.}{provider-profile}

\screenshotfig{Photos/provider/provider_account_settings.png}{Account Settings --- Password change, notification preferences, 2FA toggle, and account deletion.}{account-settings}

\screenshotfig{Photos/provider/provider_create_provider_form.png}{Create Provider Account --- First name, last name, email, and password with complexity requirements.}{create-provider}

\screenshotfig{Photos/provider/provider_create_provider_success.png}{Provider Account Created --- Confirmation toast after successful provider creation.}{provider-created}

\screenshotfig{Photos/provider/provider_messaging_center.png}{Provider Messaging Center --- Client list with real-time connection status, conversation view with message history and system notifications.}{provider-messaging}

\newpage

% ── Client Dashboard ──
\subsection{Client Dashboard (FR3, FR4, FR5, FR6, FR7, FR8, Documents)}

\screenshotfig{Photos/client/client_dashboard_overview.png}{Client Dashboard Overview --- Document count, resource count, upcoming appointments, daily mood check-in orb, suggested reading, and resource filter/search.}{client-dashboard}

\screenshotfig{Photos/client/client_mood_need_support.png}{Mood Orb: ``Need Support Now'' --- The mood orb glows red when the client needs immediate provider attention.}{mood-support}

\screenshotfig{Photos/client/client_mood_tough_day.png}{Mood Orb: ``Having a Tough Day'' --- Warm red-orange tone for the tough day state.}{mood-tough}

\screenshotfig{Photos/client/client_mood_hanging_in_there.png}{Mood Orb: ``Hanging in There'' --- Golden-amber orb for the neutral mood state.}{mood-neutral}

\screenshotfig{Photos/client/client_mood_doing_alright.png}{Mood Orb: ``Doing Alright'' --- Pink-magenta orb for the positive mood state.}{mood-alright}

\screenshotfig{Photos/client/client_mood_doing_well.png}{Mood Orb: ``Doing Well'' --- Bright green orb for the most positive mood state.}{mood-well}

\newpage

\screenshotfig{Photos/client/client_resource_filtered.png}{Resource Library (Filtered) --- Filtered by category ``Postpartum'' and difficulty ``Beginner.''}{resource-filtered}

\screenshotfig{Photos/client/client_resource_detail.png}{Resource Detail View --- ``Postpartum Daily Recovery'' content organized by week.}{resource-detail}

\screenshotfig{Photos/client/client_resource_library_grid.png}{Resource Library Grid --- Full resource grid filtered by the client's postpartum stage.}{resource-grid}

\screenshotfig{Photos/client/client_blog_post_detail.png}{Blog Post Detail (Client View) --- Published blog post with category tags, reading time, and full content.}{client-blog-detail}

\screenshotfig{Photos/client/client_blog_list.png}{Blog List (Client View) --- Authenticated blog list showing published posts.}{client-blog-list}

\newpage

\screenshotfig{Photos/client/client_appointments_calendar.png}{Appointments Calendar --- Monthly calendar with status legend and upcoming appointments list.}{client-calendar}

\screenshotfig{Photos/client/client_request_time_form.png}{Request a Specific Time --- Date picker, type selector, start/end time inputs, and optional notes.}{request-time}

\screenshotfig{Photos/client/client_request_time_filled.png}{Request Time (Completed Form) --- March 19, 2026, 7:02--9:02 PM (Virtual), ``Casual check in.''}{request-filled}

\screenshotfig{Photos/client/client_appointment_approved.png}{Appointment Approved --- Success toast confirms approval; appointment appears on calendar.}{appt-approved}

\screenshotfig{Photos/client/client_appointment_confirmed.png}{Appointment Confirmed --- Confirmed appointment in upcoming list with date, time, type, and status badge.}{appt-confirmed}

\newpage

\screenshotfig{Photos/client/client_notifications_and_chat.png}{Client Notifications and Real-time Chat --- Notification badge with unread messages and real-time chat with provider.}{client-chat}

\screenshotfig{Photos/client/client_messaging_with_notifications.png}{Client Messaging (Full View) --- Full messaging page with notification area and conversation history.}{client-messaging}

\newpage

\screenshotfig{Photos/client/client_profile_edit_and_upload.png}{Client Profile and Document Upload --- Profile edit form with sub-navigation and Upload Documents section.}{client-profile}

\screenshotfig{Photos/client/client_document_upload_form.png}{Document Upload Form --- File selector, title, description, category dropdown, and tags.}{upload-form}

\screenshotfig{Photos/client/client_document_upload_filled.png}{Document Upload (Completed Form) --- Selected image, title ``Intake,'' category ``Other,'' tag ``General.''}{upload-filled}

\screenshotfig{Photos/client/client_document_upload_success.png}{Document Upload Success --- Confirmation toast ``Document uploaded successfully!''}{upload-success}

\screenshotfig{Photos/client/client_profile_security_2fa.png}{Client Security Settings --- Two-Factor Authentication (2FA) status with manage option.}{client-2fa}

\end{document}
